% ---------------------------------------------------------------------------
% Definition of the training set
% Sources: bj_data.zip (aSTEP files), beijing_20200_20261114-1424_jesper.zip (trained data,
%          same sha256 as TSMini/data/beijing_20200.h5 and _traj_simi_dict_frechet.pkl),
%          TSMini/export_beijing.py, TSMini/prepare_rome.py, TSMini/utils/data_loader.py,
%          TrajCL/utils/data_loader.py, ConvTraj/prepare.py, ConvTraj/mynn/Traj_KNN.py,
%          branch tsmini-fraction-experiment (commits 9da3212, be002ef),
% Uses only standard LaTeX (no extra packages).
% ---------------------------------------------------------------------------
\subsection{Definition of the training set}
\label{sec:training-set}

To compare the models fairly, they must learn from the same data and be tested on
the same data. This section defines the data every model uses: where the
trajectories come from, how they are filtered and split, and how the training set
is reduced to 20\,\%, 40\,\%, 60\,\%, 80\,\% and 100\,\% of its size. All models
are trained on all five sizes; the results are presented in
Section~\ref{sec:fraction-results}.

\paragraph{The aSTEP files}
As described in Section~\ref{sec:beijing-dataset}, the Beijing data from aSTEP
consists of a training set in five sizes, a validation set and a test set
(Table~\ref{tab:astep-files}). The training files are nested: every trajectory in
\texttt{traj\_train\_20} is also in \texttt{traj\_train\_40}, and so on up to
\texttt{traj\_train\_100}. The validation and test files have no trajectories in
common with each other or with the training files.

\begin{table}[htbp]
  \centering
  \begin{tabular}{|l|r|}
    \hline
    \textbf{File} & \textbf{Trajectories} \\
    \hline
    \texttt{traj\_train\_20.csv}  & 372{,}102 \\
    \texttt{traj\_train\_40.csv}  & 744{,}204 \\
    \texttt{traj\_train\_60.csv}  & 1{,}116{,}306 \\
    \texttt{traj\_train\_80.csv}  & 1{,}488{,}408 \\
    \texttt{traj\_train\_100.csv} & 1{,}860{,}507 \\
    \texttt{traj\_valid.csv}      & 620{,}169 \\
    \texttt{traj\_test.csv}       & 620{,}169 \\
    \hline
  \end{tabular}
  \caption{The trajectory files in the Beijing dataset from aSTEP.}
  \label{tab:astep-files}
\end{table}

\paragraph{Sample used in this project}
Computing the exact Fréchet distance between all pairs of 1.86 million trajectories
is not feasible, so the models are trained on a sample. All trajectories in the
sample are taken from \texttt{traj\_train\_100}, and the smaller training sets are
created as subsets of our own training split (see \emph{Training subsets} below)
instead of from the four smaller files. This keeps the evaluation and test sets
exactly the same for every size, the Fréchet ground truth only has to be computed
once, and the smaller training sets are guaranteed to be contained in the larger
ones, just like in the aSTEP files.

The trajectories are exported from the shared aSTEP database by the script
\texttt{export\_beijing.py}. A trajectory is built from its map-matched road points
and not from its GPS points, because the GPS points usually only contain the start
and end of a trip. Each road point refers to an edge in the road network, and the
edge is replaced by its start and end nodes, so a trajectory becomes the sequence
of road network nodes it passes through. The order of the export is random but
repeatable, since the trajectories are sorted by a hash of the seed (2000) and the
trajectory id. The export stops when 50{,}000 trajectories that pass the filters
below have been found. The sample is spread evenly over the aSTEP files: for
example, 9{,}891 of the 50{,}000 trajectories (19.8\,\%) are also in
\texttt{traj\_train\_20}. None of the trajectories are from \texttt{traj\_valid} or
\texttt{traj\_test}; the evaluation and test sets used in this project are taken
from the sample.

\paragraph{Filtering and preprocessing}
A trajectory is kept only if
\begin{itemize}
  \item it lies completely inside central Beijing, between 116.20--116.55\,°E and
        39.75--40.03\,°N (about $30 \times 31$\,km, roughly inside the fifth ring
        road), and
  \item it has between 20 and 200 points.
\end{itemize}
The coordinates are then converted from longitude and latitude to Web Mercator
metres by \texttt{prepare\_rome.py}, which is used for both Rome and Beijing.
Consecutive duplicate points are removed, and the length filter is applied again.
This leaves \textbf{49{,}934 trajectories}, which are shuffled with a fixed seed
(2000) and stored as the dataset \texttt{beijing\_20200}. TrajCL and ConvTraj do not
build their own datasets but convert these files to their own formats
(\texttt{prepare\_beijing.py} and \texttt{prepare.py}), so all three models use the
same trajectories in the same order.

\paragraph{Splits}
The data loaders in TSMini and TrajCL take a fixed block of 10{,}000 trajectories
for training with a similarity measure, starting at 70\,\% of the dataset. This
block is divided into 7{,}000 training, 1{,}000 evaluation and 2{,}000 test
trajectories. ConvTraj uses the same training and test trajectories. The first
70\,\% of the dataset is only used by TrajCL's self-supervised pre-training, which
therefore never trains on the fine-tuning, evaluation or test trajectories.
Table~\ref{tab:beijing-splits} shows the parts.

\begin{table}[htbp]
  \centering
  \begin{tabular}{|l|r|r|l|}
    \hline
    \textbf{Part} & \textbf{Indices} & \textbf{Trajectories} & \textbf{Used by} \\
    \hline
    Pre-training     & 0--34{,}952        & 34{,}953 & TrajCL pre-training \\
    Training         & 34{,}953--41{,}952 & 7{,}000  & All models \\
    Evaluation       & 41{,}953--42{,}952 & 1{,}000  & Model selection \\
    Test             & 42{,}953--44{,}952 & 2{,}000  & Final scores \\
    Remaining        & 44{,}953--49{,}933 & 4{,}981  & Search only \\
    \hline
    \textbf{Total}   &                    & \textbf{49{,}934} & \\
    \hline
  \end{tabular}
  \caption{The parts of the Beijing sample (\texttt{beijing\_20200}), by index in
  the shuffled dataset. All 49{,}934 trajectories are embedded and used as the
  database when searching.}
  \label{tab:beijing-splits}
\end{table}

\paragraph{Ground truth}
The models learn to approximate the discrete Fréchet distance, computed with
Euclidean distances in the Mercator plane. The exact distances are computed once by
\texttt{prepare\_rome.py} between all pairs inside each split: $7{,}000 \times
7{,}000$ for training, $1{,}000 \times 1{,}000$ for evaluation and $2{,}000 \times
2{,}000$ for testing, about 27 million distances in total. Only the upper triangle
of each matrix is stored, and the distances are divided by the largest distance in
the ground truth before training, so all values lie between 0 and 1. Because a
training subset uses the top-left block of the training matrix, no new distances
have to be computed for the smaller training sets.

\paragraph{Training subsets}
The 20--80\,\% training sets are the first 20\,\%, 40\,\%, \ldots{} of the 7{,}000
training trajectories (Table~\ref{tab:training-subsets}). Since the dataset was
shuffled beforehand, each subset is a random sample, and the subsets are nested:
the 20\,\% set is part of the 40\,\% set, and so on. The subsets are the same for
all three models, and the evaluation and test sets are the same 1{,}000 and
2{,}000 trajectories for every size, so the only thing that changes is how many
trajectories a model learns from.

\begin{table}[htbp]
  \centering
  \begin{tabular}{|r|r|r|}
    \hline
    \textbf{Share} & \textbf{Training trajectories} & \textbf{Ground-truth pairs} \\
    \hline
    20\,\%  & 1{,}400 & 979{,}300 \\
    40\,\%  & 2{,}800 & 3{,}918{,}600 \\
    60\,\%  & 4{,}200 & 8{,}817{,}900 \\
    80\,\%  & 5{,}600 & 15{,}677{,}200 \\
    100\,\% & 7{,}000 & 24{,}496{,}500 \\
    \hline
  \end{tabular}
  \caption{The five training set sizes. Ground-truth pairs are the distinct pairs
  of training trajectories with a known Fréchet distance, $n(n-1)/2$.}
  \label{tab:training-subsets}
\end{table}

\paragraph{How the subsets are used by each model}
\begin{itemize}
  \item \textbf{TSMini.} The share is set with the option
        \texttt{trajsimi\_train\_frac}, and the script
        \texttt{run\_fraction\_experiment.py} trains all five sizes one after the
        other. Each batch consists of 128 trajectories drawn at random from the
        subset. The number of batches per epoch is calculated from the full
        training set ($\lceil (7{,}000/128)^2 \rceil = 2{,}991$) for every size. A
        smaller training set therefore gets the same number of training steps, only
        with less data, so the comparison is not affected by shorter epochs.
  \item \textbf{TrajCL.} Only the fine-tuning step learns from the 7{,}000
        training trajectories, so only this step uses the subsets. The cell
        embeddings and the pre-training use the 34{,}953 pre-training trajectories
        and are the same for every size.
        % TODO: how the share is set, and whether batches per epoch stay the same
  \item \textbf{ConvTraj.} The share is set with the option
        \texttt{train\_ratio}, which uses the first share of the training
        trajectories and therefore gives the same subsets as for TSMini.
        % TODO: note that one epoch goes through the subset once (fewer batches for smaller sets)
\end{itemize}

\paragraph{Limitations of the definition}
Since the subsets are nested and always taken from the start of the training
split, which trajectories end up in the smaller sets is decided by the shuffle. All
sizes are also taken from one sample of 49{,}934 trajectories, so the five sizes
vary the number of training trajectories with a ground truth, not the use of the
full aSTEP files.
